\documentclass[a4paper,12pt]{article}

\usepackage[utf8]{inputenc}
\usepackage[russian]{babel}
\usepackage{amsmath}
\usepackage{amssymb}
\usepackage{graphicx}

\title{Лекция 2 \\ Свойства линейных кодов и коды Хэмминга}
\date{}

\begin{document}

\maketitle

\section{Минимальное расстояние линейного кода}

Для линейного кода минимальное расстояние $d_{min}$ можно определить,
используя его проверочную матрицу $H$.

Известно, что любое кодовое слово $c$ удовлетворяет условию

\[
Hc^T = 0.
\]

Если существует ненулевое кодовое слово веса $w$, то это означает,
что некоторая линейная комбинация $w$ столбцов матрицы $H$
дает нулевой вектор.

Следовательно, минимальное расстояние кода совпадает
с наименьшим весом ненулевого кодового слова.
В терминах проверочной матрицы это означает,
что $d_{min}$ равно минимальному числу столбцов $H$,
которые образуют линейно зависимую систему.

Таким образом, задача нахождения минимального расстояния
сводится к поиску минимального количества линейно зависимых
столбцов в проверочной матрице.

\section{Критерий минимального расстояния}

Сформулируем важное утверждение.

\textbf{Теорема.}
Пусть задан линейный $(n,k)$-код с проверочной матрицей $H$.
Минимальное расстояние кода равно $d$ тогда и только тогда,
когда выполняются два условия:

\begin{itemize}
\item любые $d-1$ столбцов матрицы $H$ линейно независимы;
\item существует набор из $d$ столбцов, которые являются
линейно зависимыми.
\end{itemize}

Иначе говоря, $d$ является минимальным количеством столбцов
проверочной матрицы, которые образуют линейно зависимую систему.

Это свойство широко используется как при анализе кодов,
так и при построении кодов с заданной корректирующей способностью.

\section{Дуальные коды}

В теории кодирования важным понятием является \textbf{дуальный код}.

Пусть $C$ — линейный $(n,k)$-код над полем $F_2$.
Тогда дуальным кодом $C^\perp$ называется множество всех
векторов длины $n$, которые ортогональны каждому
кодовым словам из $C$:

\[
C^\perp = \{ v \in F_2^n \mid c \cdot v = 0 \ \text{для всех } c \in C \}.
\]

Размерность дуального кода равна $n-k$.

Порождающей матрицей дуального кода служит
проверочная матрица исходного кода.

\subsection{Пример}

Рассмотрим код Хэмминга $(7,4)$.
Его дуальный код имеет параметры $(7,3)$
и называется \textbf{симплекс-кодом}.

Особенностью симплекс-кода является то,
что все его ненулевые кодовые слова
имеют одинаковый вес, равный 4.

\section{Оптимальность кодов Хэмминга}

Двоичные коды Хэмминга обладают важным свойством —
они являются оптимальными для своего набора параметров.

Пусть длина кода равна

\[
n = 2^r - 1
\]

и минимальное расстояние равно $3$.

Максимально возможное количество кодовых слов
в любом коде с такими параметрами ограничено
\textbf{границей Хэмминга}:

\[
M \le \frac{2^n}{1+n}.
\]

Коды Хэмминга с параметрами

\[
(2^r - 1,\; 2^r - r - 1)
\]

достигают этой границы. Это означает,
что они содержат максимально возможное
число кодовых слов.

Поэтому такие коды называют
\textbf{совершенными кодами}.

Другими словами, не существует кода
с теми же параметрами $n$ и $d=3$,
который содержал бы больше кодовых слов.

\section{Симплекс-коды}

Дуальные коды к кодам Хэмминга образуют
класс кодов, называемых \textbf{симплекс-кодами}.

Если код Хэмминга имеет параметры

\[
(2^r - 1,\; 2^r - r - 1),
\]

то соответствующий симплекс-код имеет параметры

\[
(2^r - 1,\; r).
\]

Главная особенность симплекс-кодов состоит в том,
что все ненулевые кодовые слова имеют одинаковый вес:

\[
2^{r-1}.
\]

\subsection{Пример}

При $r = 3$ код Хэмминга имеет параметры $(7,4)$.
Дуальный ему симплекс-код имеет параметры $(7,3)$.

Он содержит $8$ кодовых слов,
и каждое ненулевое слово имеет вес $4$.

\section{Расширенные коды Хэмминга}

Расширенный код Хэмминга получается из обычного
кода Хэмминга путем добавления одного
дополнительного бита четности.

Если исходный код имеет параметры

\[
(2^r - 1,\; 2^r - r - 1,\; 3),
\]

то после добавления общего бита четности
получается код с параметрами

\[
(2^r,\; 2^r - r - 1,\; 4).
\]

Дополнительный символ равен сумме всех
битов кодового слова по модулю два.

Добавление этого бита увеличивает
минимальное расстояние до $4$.

Это позволяет:

\begin{itemize}
\item исправлять одну ошибку,
\item обнаруживать две ошибки.
\end{itemize}

\subsection{Пример}

Хорошо известным примером является
расширенный код Хэмминга

\[
(8,4,4).
\]

Он широко применяется в системах
обнаружения и исправления ошибок.

Интересно, что этот код является
\textbf{самодуальным}, то есть совпадает
со своим дуальным кодом.

\end{document}